\noindent
\graphicspath{{./22-08/Images/}} 
%%%%%%%%%%%%%%%%
% Page de gauche
%%%%%%%%%%%%%%%%
\begin{minipage}[t][\textheight][t]{0.33\textwidth}
    \colorbox{StrongBlue}{
        \parbox[c][\textheight][c]{\textwidth}{ 
            \color{OffWhite} 
            \begin{center}
                \begin{minipage}[t][0.6\textheight][c]{0.8\textwidth}
                    {\huge \textbf{Vendredi 22 Août}}\\

                    \normalsize
                    \begin{center}
                        \begin{tabular}{@{} l @{\hspace{0.5cm}} l @{}}
                            Abri de départ : &\hfill Port Haliguen \\ 
                            Abri d'arrivée : &\hfill Penboch \\ 
                            Distance : &\hfill 18 Milles \\ 
                            Départ : &\hfill 09 h 00 \\ 
                            Arrivée prévue : &\hfill 12 h 00 \\ 
                            Lever soleil 
\includegraphics[width=8px,height=8px]{../../Pictos/soleil_OffWhite.png} : &\hfill 07 h 17 \\ 
                            Coucher soleil 
\includegraphics[width=8px,height=8px]{../../Pictos/lune_OffWhite.png} : &\hfill 21 h 09 \\ 
                            \\
                            Abris replis : &  \\
                                    & • Port Haliguen \\
                                    & • Crouesty \\
                                    & • Port Navalo \\
                                    & • Anse de Penhap \\                    
                                
                        \end{tabular}
                    \end{center}
                    
                \end{minipage}%

                \begin{minipage}[t][0.4\textheight][t]{\textwidth}
                    \begin{center}
    {\large \textit{Marées \textbf{Port Haliguen}}}\\ \vspace{0.5em}
    
    \renewcommand{\arraystretch}{1.40} % Adjust the row height by a factor of 1.4
    
       \begin{tabular}{ |C{5em}|c|C{4em}|C{4em}| } \hline
        & Coef & Heure & Hauteur\\ \hline
        PM & \multirow{2}{1em}{78}  & 05 h 18 & 4.84 m\\ \cline{1-1} \cline{3-4} 
        BM &                        & 11 h 12 & 1.08 m\\ \hline
        PM & \multirow{1}{1em}{83}  & 17 h 30 & 5.05 m\\ \hline
    \end{tabular}
    
\end{center}

                    \begin{center}
    {\large \textit{Marées \textbf{Port Navalo}}}\\ \vspace{0.5em}
    
    \renewcommand{\arraystretch}{1.40} % Adjust the row height by a factor of 1.5
    \begin{tabular}{ |C{5em}|c|C{4em}|C{4em}| } \hline
        & Coef & Heure & Hauteur\\ \hline
        BM & \multirow{1}{1em}{78}  & 11 h 16 & 1.06 m\\ \hline
        PM & \multirow{2}{1em}{83}  & 17 h 42 & 4.90 m\\ \cline{1-1} \cline{3-4}
        BM &   				& 23 h 46 & 0.82 m\\ \hline
    \end{tabular}
\end{center}
                \end{minipage}%
            \end{center}       
        }      
    }
\end{minipage}%
\begin{minipage}[c][\textheight][t]{0.10\textwidth}
\hfill
\end{minipage}%
\begin{minipage}[c][\textheight][t]{0.50\textwidth}
    \vspace{3cm}
    \color{StrongBlue}    
    {\huge \textbf{Description de la navigation}}\\
            \begin{itemize}
                \item Journée de retour dans le Golfe et plus grosse journée de navigation du camp. 
                \item Départ à 9h, cap 79° depuis la cardinale \textit{Quiberon Sud}. On attend d'avoir largement laissé la cardinale \textit{Sud Méhaban} avant de mettre le cap plein Nord vers la jument. Si le virement est déclenché trop tôt la flottille peut arriver sur \textit{Baguen Hir}, on s'assurera donc de laisser la cardinale largement à notre Ouest. 
                \item Arrivée devant la jument un peu avant midi, si les conditions de vent le permette on pourra envisager de passer à la voile. Dans tous les cas on sort la carte dans le cockpit et on allumera le moteur au point mort par sécurité. La consigne sera donnée d'enrouler le gênois et de faire route moteur en cas de difficulté. 
                \item Déjeuner au mouillage pointe de Kerners, départ vers 14h pour Penhap en passant au Sud de l'Île aux moines, pour varier les itinéraires et maximiser le temps de navigation. Si besoin, la navigation peut être écourtée en passant par le Nord.
                \item Choix de coffres le plus groupés possible pour le mouillage. 
                \item Veillée d'envoi des troisièmes années sur le terrain du centre spirituel de Penboch. 
            \end{itemize}
        {
        \huge \textbf{Contraintes de la navigation}}\\
            \begin{itemize}
                \item Passage du courant de la Jument en vives eaux, il est nécessaire de passer avec la marée montante, idéalement en évitant les 3dz : passage donc entre 11h15 et 13h15. 
                \item Beaucoup de balisage et d'obstacles dans la jument, sans tous les énumérer on aura une attention toute particulière à ce que les bateaux restent à distance raisonnable de la marque latérale tribord du \textit{Grand Mouton} et de l'îlot \textit{Er Lannig} dont la côte n'est pas franche. 
                \item Le Golfe abrite des effets de sites pour le vent et le courant, on sera donc vigilant à l'évolution des conditions.  
                \item Le passage entre Arz et l'Île aux moines n'est pas très franc, on veillera à respecter le balisage dans le sens conventionelle d'entrée au port. 
            \end{itemize}
\end{minipage}%
%%%%%%%%%%%%%%%%
% Page de droite
%%%%%%%%%%%%%%%%
\newgeometry{top=2cm, bottom=2cm, left=2cm, right=2cm}
\begin{minipage}[t][0.9\textheight][t]{0.60\textwidth}
    \color{StrongBlue}    
    \noindent
    \begin{minipage}[t][0.3\textheight][t]{\textwidth}
        \begin{minipage}[t][0.3\textheight][t]{0.65\textwidth}
            \vspace{-8.5em}
            {\huge \textbf{Détails de la navigation}}\\
            Le QR-code situé à droite est cliquable, il permet d'explorer le planning de navigation sur le site datashom. 
            La carte ci-dessous visualise le trajet prévu ainsi que les abris de replis et dangers identifiés. 
        \end{minipage}%
        \hfill
        \begin{minipage}[t][0.3\textheight][t]{0.25\textwidth}       
        \href{https://data.shom.fr/datacontext/fr/1748012047608-918142}{ 
            \includegraphics[width=0.8\textwidth]{qr-code.png}
        }
        \end{minipage}%
    \end{minipage}
    \begin{minipage}[t][0.7\textheight][t]{\textwidth}
        % \includegraphics[width=0.97\textwidth]{carte.png}
        \includegraphics[height=0.7\textheight]{carte.png}
    \end{minipage}
\end{minipage}
\begin{minipage}[t][0.9\textheight][t]{0.40\textwidth}
    \color{StrongBlue} 
    \centering 
    \includegraphics[height=\textheight, width=\textwidth, keepaspectratio]{illustration_verticale.jpg}
\end{minipage}%
\restoregeometry



